% !PRINTED 527
\weberpage{527}
Ebenso können nicht 1 und 3 cyklisch an einander stossen, und es bleiben also nur folgende übrig:
\weq{4}{(0,1,2,3), \quad (0,3,2,1), \quad (0,3)(1,2), \quad (0,1)(2,3).}
Wenn aber von diesen Permutationen eine zu (3) hinzutritt, so müssen auch alle anderen aufgenommen werden, denn es ist
\[
\begin{gathered}
(0,2)(0,1,2,3) = (0,3)(1,2) \\
(1,3)(0,1,2,3) = (0,1)(2,3) \\
(0,2)(1,3)(0,1,2,3) = (0,3,2,1),
\end{gathered}
\]
und man kann also durch Vermittelung der Transpositionen $(0,2)$ und $(1,3)$ jede der Permutationen (4) aus jeder anderen herleiten. Aus der Voraussetzung also, dass überhaupt in einer von $P$ und $Q$ verschiedenen transitiven Gruppe eine Transposition vorkommt, folgt, dass dies die Gruppe achten Grades
\weq{5}{\begin{gathered}
P_1 = 1, \quad (0,2), \quad (1,3), \quad (0,2)(1,3), \quad (0,1)(2,3), \quad (0,3)(1,2) \\
(0,1,2,3), \quad (0,3,2,1)
\end{gathered}}
oder wenigstens eine damit conjugirte sein muss.

Dass $P_1$ wirklich eine Gruppe ist, erkennt man leicht, wenn man je zwei seiner Permutationen zusammensetzt.

Es giebt drei verschiedene mit $P_1$ conjugirte Gruppen, die man erhält, wenn man statt von $(0,2)$, $(1,3)$ von den Transpositionen $(0,1)$, $(2,3)$ oder $(0,3)$, $(1,2)$ ausgeht. Der Durchschnitt dieser drei Gruppen ist ein Normaltheiler von $P$ und zugleich von $Q$, nämlich
\weq{6}{Q_1 = 1, \quad (0,2)(1,3), \quad (0,1)(2,3), \quad (0,3)(1,2).}
Wenn eine transitive, von $P$ und $Q$ verschiedene Gruppe keine Transposition enthält, so enthält sie entweder keinen viergliedrigen Cyklus und ist dann mit (6) identisch, oder sie enthält einen viergliedrigen Cyklus und enthält dann auch die ganze cyklische Gruppe, wie
\weq{7}{P_2 = 1, \quad (0,1,2,3), \quad (0,2)(1,3), \quad (0,3,2,1).}
Enthält eine solche Gruppe noch einen zweiten viergliedrigen Cyklus, etwa $(0,1,3,2)$, so ist sie wegen $(0,1,2,3)(0,1,3,2) = (0,3,1)$ die symmetrische Gruppe. Enthält sie aber ausser $P_2$ noch eine der Permutationen von (6), etwa $(0,1)(2,3)$, so ist sie wegen
\[
(0,1,2,3)(0,1)(2,3) = (1,3)
\]
mit $P_1$ identisch.
